\section{From associative equivalence to Poisson equivalence}
\label{sec:formality}

We now work over $\CC$, using the embedding of the coefficient field
chosen above. We write $h$ for the deformation parameter and $t$
for the Rees parameter. In this section $W$ denotes the complex scalar
extension of the common vector space of the two Rees families.
Put $A=S_{\CC}(W)$ and write $\pi_L(t),\pi_M(t)$ for their linear
Poisson tensors.  Their common specialization is denoted by $\pi_0$.
For a complex vector space $X$, our conventions are
\[
 X((h))=X[[h]][h^{-1}],\qquad
 E=\CC((h)),\qquad R=E[[t]],\qquad B=A((h)).
\]
The order is completion in $h$, localization, then completion in
$t$.  Each $t$-coefficient has its own lower Laurent bound.  We do not
require a bound independent of the $t$-degree.  In particular,
When $W\ne0$, $B$ strictly contains $E[W]$; for $W=0$ both equal $E$.

\begin{proposition}\label{prop:poisson-comparison}
The marked Rees isomorphism of \cref{prop:rees-isomorphism} induces an
$R$-algebra automorphism $\Psi$ of $B[[t]]$, congruent to the identity
modulo $t$, such that
\[
 \Psi\bigl(\{f,g\}_{\pi_L(t)}\bigr)
   =\{\Psi(f),\Psi(g)\}_{\pi_M(t)}
 \qquad(f,g\in B[[t]]).
\]
\end{proposition}

The proposition concerns the indicated completed function algebra,
not merely finite polynomial inputs or unrelated finite-order
approximations.  We prove it using affine formality and a complete
filtered comparison theorem.

\subsection*{Full cochains and their completed action}

Use the shifted differential graded Lie algebras
\[
 P_{\mathrm{poly}}^p=A\otimes\Lambda^{p+1}W^*,\quad -1\le p\le\dim W-1,
 \qquad
 C_{\mathrm{all}}^p=\Hom_{\CC}(A^{\otimes(p+1)},A),\quad p\ge-1.
\]
All completions of these graded complexes are taken degreewise.
The first has zero differential and the Schouten bracket; the second
has the Gerstenhaber bracket and differential $[m,-]$, where $m$ is
ordinary multiplication.  Both contain functions in degree $-1$ and
vanish in lower degrees.  In degree zero, $C_{\mathrm{all}}^0=\End_{\CC}(A)$
with its commutator bracket.  We normalize bivectors so that the first
star commutator is $h$ times their Poisson bracket.

The factorial-normalized Hochschild--Kostant--Rosenberg map sends a
wedge of $q$ derivations to $1/q!$ times their alternating multilinear
product.  It is a quasi-isomorphism into the \emph{full} Hochschild
complex \cite{HKR1962}.  In the polynomial case this follows directly
from the Koszul resolution of the diagonal.  Indeed, the elements
$x_i\otimes1-1\otimes x_i$ in $A^e$ become independent polynomial
variables after a change of coordinates, hence their Koszul complex
resolves $A$ \cite{StacksKoszul}.  Applying $\Hom_{A^e}(-,A)$ gives zero
differentials and the polyvector spaces.  The alternating comparison
with the free bar resolution computes full Hochschild cohomology;
its cochain composite with normalized HKR is the identity.  The degree
signs between the bar differential and $[m,-]$ are absorbed into this
comparison.

Kontsevich's affine formality theorem supplies an $L_\infty$ morphism
$\mathcal U$ from polynomial polyvectors to polynomial polydifferential
cochains, with first component HKR
\cite[\S\S4.6.2, 6.4 of the arXiv version]{Kontsevich2003}.
We fix the morphism given by his affine graph construction. Its formulas preserve
polynomials, and their real coefficients extend to $\CC$.  Composing
with the strict inclusion into $C_{\mathrm{all}}$ gives an $L_\infty$
quasi-isomorphism with the full target just described.

Here are the elementary convergence facts we will use.  An element
\[
 F=\sum_{j\ge a}h^jF_j
 \in\Hom_{\CC}(A^{\otimes q},A)((h))
\]
acts $E$-multilinearly on $B^q$.  On inputs with lower bounds $b_i$,
the coefficient of $h^d$ sums over
$j+j_1+\cdots+j_q=d$ with $j\ge a$ and $j_i\ge b_i$, a finite set.
The same argument proves compatibility with compositions and brackets.
The action is faithful, by evaluating on constant polynomial inputs.
It is not an identification with all $E$-linear maps on $B$.
After adjoining $t$, first split a fixed $t$-degree in finitely many
ways, and then apply this Laurent argument.  Consequently identities
between these completed operators extend from polynomial inputs by
coefficientwise expansion, without any assumption that arbitrary
linear maps are continuous.

In particular, a Laurent series of unary cochains has a lower bound
for the whole operator, rather than a bound chosen separately after
seeing its input.  If two such operators have bounds $a,b$, their
composition has bound $a+b$.  For an outer series
$F=\sum_{r\ge0}t^rF_r$ acting on
$f=\sum_{s\ge0}t^sf_s$, the coefficient of $t^n$ is
$\sum_{r+s=n}F_r(f_s)$.  Its Laurent lower bound is the minimum of
finitely many sums of bounds, hence is finite.  This explains why
the bounds may vary with $r$, and why no exchange of the two
completions is implicit.  The same reasoning applies to every finite
arity occurring in the deformation complexes.

\subsection*{Twisting before localization}

Set $\alpha_0=h\pi_0$ and
\[
 \mu_{\star,0}=m+\sum_{j\ge1}\frac1{j!}\mathcal U_j(\alpha_0^{\otimes j}).
\]
The twisted components are
\[
 \mathcal U_n^{\alpha_0}(z_1,\ldots,z_n)
 =\sum_{j\ge0}\frac1{j!}
   \mathcal U_{n+j}(\alpha_0^{\otimes j},z_1,\ldots,z_n).
\]
Every additional base input increases $h$-order by one.  These sums
therefore converge before localization, and the twisted differentials
are $[\alpha_0,-]$ and $[\mu_{\star,0},-]$.  The twisted identities follow by
comparing finite coefficients in the original identities.

\begin{lemma}\label{lem:twisted-quasi-iso}
The twisted linear component induces a quasi-isomorphism
\[
 \bigl(P_{\mathrm{poly}}((h)),[\alpha_0,-]\bigr)
 \longrightarrow
 \bigl(C_{\mathrm{all}}((h)),[\mu_{\star,0},-]\bigr).
\]
\end{lemma}
\begin{proof}
First form its cone over $\CC[[h]]$. Write $C$ for the underlying
graded vector space of the ordinary HKR cone. Its underlying module
is $C[[h]]$, with differential $d=d_0+O(h)$, and its reduction
modulo $h$ is the acyclic HKR cone.  Choose a contracting homotopy
$s_0$ for this vector-space complex and extend it coefficientwise.
Then
\[
 \mathcal K=ds_0+s_0d=1+O(h)
\]
is invertible by a geometric series.  Since $d^2=0$, $\mathcal K$ and $\mathcal K^{-1}$
commute with $d$.  Thus $s=s_0\mathcal K^{-1}$ satisfies $ds+sd=1$.  Every
coefficient is finite: $\mathcal K-1$ increases $h$-order.  No polynomial
degree bound on $s_0$ is needed.  Localizing this contraction gives a
contraction of the Laurent cone, proving the assertion.
\end{proof}

The higher twisted components also extend to Laurent inputs: a fixed
finite list of such inputs has a finite sum of lower bounds, while
the number of base insertions still increases $h$-order.  Notice
that quasi-isomorphism was proved for the completed cone \emph{before}
inverting $h$.

\subsection*{Complete filtered gauge reflection}

Over $E$, form
\[
 D_P=t\bigl(P_{\mathrm{poly}}((h)),[\alpha_0,-]\bigr)[[t]],
 \qquad
 D_C=t\bigl(C_{\mathrm{all}}((h)),[\mu_{\star,0},-]\bigr)[[t]],
\]
with $F^r=t^r(-)[[t]]$, $r\ge1$.  These filtrations are separated and
complete, start with the whole space, and satisfy
$[F^r,F^s]\subset F^{r+s}$.  The twisted Taylor maps preserve sums of
filtration degrees.  At a fixed $t$-degree $r$, their MC series uses
at most $r$ non-base inputs and finitely many partitions of $r$;
the remaining Laurent sums are already justified.

On every $F^r$, the linear component is the product of copies of the
map in \cref{lem:twisted-quasi-iso}, indexed by $t$-degrees at least $r$.
Its cone is contractible by applying the preceding contraction to
each coefficient.  Hence it is a quasi-isomorphism on \emph{every}
filtration step.  All hypotheses of the complete filtered
Dolgushev--Rogers theorem are now satisfied over the characteristic-zero
field $E$ \cite[Theorem~2.2]{DolgushevRogers2015}.  We obtain a bijection
on connected components of the completed MC simplicial sets.
The theorem is applied to these $t$-positive objects, not directly to
all Laurent polyvectors: positivity provides pronilpotence even
though negative powers of $h$ occur.  The differentials are
independent of $t$, so the coefficientwise contraction really
preserves each $F^r$.  The shifted convention used in the cited
theorem is obtained by shifting our complexes by one; it imposes no
extra nonnegativity condition on our cochain degrees.

For these DGLAs the components are gauge orbits under
$\exp(F^1D^0)$.  Indeed, an MC one-simplex has the form
$\beta(s)+ds\,\xi(s)$; its equation is the gauge differential equation.
Modulo $t^r$, the time-ordered solution with identity initial value
involves finitely many
iterated integrals of polynomials in $s$, defined over any
characteristic-zero field. Its uniqueness makes these solutions compatible
as $r$ increases, and their limit gives one completed gauge element.
A finite chain of one-simplices gives a product of these elements.
Conversely a gauge
exponential gives an MC path.  Degree $-1$ does not add a component
to a one-simplex, because the interval has no nonzero two-forms.
Thus we reflect genuine completed gauge equivalence, rather than
choose unrelated equivalences in all finite quotients.

For a differential $[a_0,-]$, the gauge action is characterized by
\[
 a_0+e^X\mathbin{\cdot}\beta
   =e^{\ad X}(a_0+\beta).
\]
Equivalently the action on $\beta$ is
$e^{\ad X}\beta-\sum_{j\ge0}(\ad X)^j([a_0,X])/(j+1)!$.
We use this with $a_0=\alpha_0$ and $a_0=\mu_{\star,0}$, respectively.

\subsection*{The actual associative comparison}

Let $\star_L(t),\star_M(t)$ be the star products obtained from the
same $\mathcal U$.  For a linear Poisson tensor, Kontsevich's
linear-Poisson comparison gives
\[
 x\star y-y\star x=h[x,y]_t\qquad(x,y\in W),
\]
and its top polynomial-degree product is ordinary multiplication
\cite[\S8.3.1 of the arXiv version]{Kontsevich2003}.  Universality and PBW triangularity
therefore give an algebra isomorphism
\[
 I_L:U_h(L_t)\longrightarrow(A[t,h],\star_L),
 \qquad I_L(x)=x,
\]
and similarly $I_M$.

These are compatible family maps.  Their values on monomials are
products of the corresponding linear generators.  A graph term of
order $h^j$ for a linear Poisson tensor contributes $j$ linear
coefficients and $2j$ derivatives, so it lowers total polynomial
degree by $j$.  It vanishes when that resulting degree is negative.
Thus each polynomial calculation is finite and polynomial in the
structure constants, and works over $\CC[t,h]$.  Define
\[
 J_L=I_L\circ\operatorname{sym}_L,\qquad
 J_M=I_M\circ\operatorname{sym}_M.
\]
They intertwine symmetric PBW products with star products, fix linear
functions, and have the form $1+O(h)$.  Their inverses exist by
the $h$-adic geometric series, with no negative $h$-powers.
Consequently all four operators belong to
$\End_{\CC}(A)[[h]][[t]]$.  Since the two brackets agree at zero,
the same construction gives $J_L(0)=J_M(0)=J_0$.

More explicitly, $J_L(x_{i_1}\cdots x_{i_q})$ is the average over
permutations of the $q$ corresponding star factors.  Thus its value
at the common special bracket is forced; no pointwise choice of an
isomorphism remains.  Triangularity proves bijectivity by induction
on polynomial degree, and the same induction works with polynomial
$t$-coefficients.  Passing to series then uses the already defined
coefficientwise action.  Ordinary PBW symmetrization is only a linear
identification here.  It is $J_L$, and respectively $J_M$, that
converts its transported product into the chosen star product.

Write $T$ for the completed homogenized Rees comparison supplied by
\cref{prop:Laurent-operator-bound}.  Its essential conclusion is
\[
 T\in1+t\End_{\CC}(A)((h))[[t]],\qquad
 \operatorname{ord}_{h}(T_r)\ge-r(D-1),
\]
with the bound valid simultaneously on all polynomial inputs.
Here $T_r$ is the coefficient of $t^r$, and $D$ is the finite
ordinary PBW degree bound from the preceding section.  In symmetric
coordinates a block from degree $a$ to degree $b$ is multiplied by
$h^{a-b}$; the bound $b\le a+r(D-1)$ is precisely what gives the
displayed estimate uniformly over all $a$.  This identification
retains the actual original enveloping map, rather than introducing
an abstract equivalence between deformation classes.  It intertwines
the actual completed PBW products.  The identity on
finite polynomial inputs extends to $B[[t]]$ by the finite-convolution
argument above; no globally defined inverse degree-scaling map on
$B$ is asserted.

Now set
\[
 G=J_M T J_L^{-1}.
\]
This is an isomorphism from $\star_L$ to $\star_M$ on $B[[t]]$, and
$G(0)=J_0\Id J_0^{-1}=\Id$.  For each $t$-degree, its coefficients
are finite sums of compositions of uniformly bounded Laurent
operators with nonnegative $h$-series.  Therefore
\[
 X=\log G\in t\End_{\CC}(A)((h))[[t]]=D_C^0.
\]
At outer degree $r$, the logarithm uses at most $r$ positive factors;
this proves its existence in the stated operator algebra.  Its
coefficients need not be polydifferential, which is why the full
Hochschild target was necessary.

If $m_L,m_M$ denote the completed star products, multiplicativity
says
\[
 m_M=G\circ m_L\circ(G^{-1}\otimes G^{-1})
      =e^{\ad X}m_L.
\]
The second identity follows by differentiating conjugation by
$e^{sX}$; its infinitesimal action is the Gerstenhaber bracket.
This supplies a target gauge equivalence relative to the common $\mu_{\star,0}$.

\begin{proof}[Proof of \cref{prop:poisson-comparison}]
Set $\beta_L=h(\pi_L(t)-\pi_0)$ and
$\beta_M=h(\pi_M(t)-\pi_0)$.  They are MC elements of
$D_P$.  Expansion of the twisted Taylor sums gives exactly
\[
 (\mathcal U^{\alpha_0})_*(\beta_L)=m_L-\mu_{\star,0},
 \qquad
 (\mathcal U^{\alpha_0})_*(\beta_M)=m_M-\mu_{\star,0}.
\]
Grouping terms by their base and positive-$t$ inputs produces the
factors $1/j!$ and $1/n!$; all rearrangements are justified by the
two coefficientwise convergence arguments.  Thus the constructed
$G$ compares precisely these two MC images.  Injectivity on connected
components reflects this equivalence to the source.

A source gauge element is the exponential of a $t$-positive vector
field with coefficients in $B[[t]]$.  Its action on functions is
well-defined by Laurent convolution and $t$-adic convergence.
The coefficientwise Leibniz rule implies that its exponential is
an $R$-algebra automorphism, with inverse the opposite exponential,
and it is the identity modulo $t$.  The source gauge equation gives
the required Poisson intertwining for $h\pi_L(t)$ and
$h\pi_M(t)$, with the direction reversed if necessary by inversion.
Cancel the invertible scalar $h$ to obtain the assertion.
\end{proof}
